% Source: https://artofproblemsolving.com/community/c3h1976603_geo_problem
% Stable identity: aops-topic-1976603-q01
Let $K$ be the foot of the perpendicular from $J$ to $GH$, and let $h=JK$. Segments $FG$, $JK$ and $IH$ are all perpendicular to $GH$, hence parallel.

\begin{asy}
usepackage("amsmath");
size(4.5cm);
defaultpen(fontsize(10pt));
pair G=(0,0),H=(6,0),F=(0,4),I=(6,8);
pair J=extension(F,H,G,I);
pair K=(J.x,0);
draw(G--H--I--cycle,linewidth(1));draw(G--F--H,linewidth(1));
draw(J--K,dashed);
real s=0.35;
draw((s,0)--(s,s)--(0,s));draw((6-s,0)--(6-s,s)--(6,s));draw(K+(s,0)--K+(s,s)--K+(0,s));
dot(J);
label("$G$",G,SW);label("$H$",H,SE);label("$F$",F,W);label("$I$",I,E);label("$J$",J,N);label("$K$",K,S);
label("$4$",(0,2),W);label("$8$",(6,4),E);label("$h$",(J+K)/2,W);
label("\mbox{}",(min(currentpicture,user=true)+max(currentpicture,user=true))/2,black); // keeps the text after the figure black when LaTeX typesets the labels inline
\end{asy}

Triangles $HKJ$ and $HGF$ are similar, so $\dfrac{h}{4}=\dfrac{KH}{GH}$. Triangles $GKJ$ and $GHI$ are similar, so $\dfrac{h}{8}=\dfrac{GK}{GH}$. Adding, since $GK+KH=GH$,
\[
\dfrac {h}{4}+\dfrac {h}{8}=1\quad\Longrightarrow\quad\dfrac{3h}{8}=1\quad\Longrightarrow\quad h=\boxed{\dfrac{8}{3}}.
\]
The length $GH$ does not matter. (In general $h=\dfrac{ab}{a+b}$, half the harmonic mean of the two heights.)
